\CNsection{1}{5G Core Overview}

The 5G Core (5GC) network is defined in \textbf{3GPP TS 23.501} and represents a fundamental redesign of the mobile core network. Unlike the 4G EPC which used point-to-point interfaces between monolithic network elements, the 5GC adopts a \textbf{Service-Based Architecture (SBA)} where Network Functions (NFs) expose services via RESTful APIs over HTTP/\allowbreak{}2.

\CNsubsection{}{Key Design Principles}

\begin{itemize}
\tightlist
\item
  \textbf{Service-Based Architecture (SBA):} NFs communicate via service-based interfaces (SBI) using HTTP/\allowbreak{}2 and JSON
\item
  \textbf{Control and User Plane Separation (CUPS):} Complete decoupling of CP and UP
\item
  \textbf{Stateless NF Design:} State stored externally (UDSF) enabling horizontal scaling
\item
  \textbf{Network Slicing Native:} Built-in support for multiple logical networks on shared infrastructure
\item
  \textbf{Cloud-Native:} Designed for containerized, microservice-based deployment
\item
  \textbf{Registration-Based NF Discovery:} NFs register with NRF and discover peers dynamically
\end{itemize}

\CNkeepfig{m19-sbi-view}{12}

\CNsection{2}{Reference Architecture}

The 5GC architecture has two equivalent representations defined in TS 23.501:

\CNsubsection{2.1}{Service-Based Interface (SBI) Representation}

In this view, all control-plane NFs connect to a common service bus and expose/\allowbreak{}consume services.

\CNfigure{m19-sbi-view}{5G core in the service-based representation}

\CNkeepfig{m19-reference-points}{5}

\CNsubsection{2.2}{Reference Point Representation}

In this view, explicit point-to-point reference points are shown between NFs.

\CNfigure{m19-reference-points}{5G core in the reference-point representation}

\Needspace{14\baselineskip}

\CNsection{3}{Network Functions in Detail}

\CNchained\CNsubsection{3.1}{Access and Mobility Management Function (AMF)}

The AMF is the single point of contact for all UE-related signaling from the RAN. It terminates the N1 (NAS) and N2 (NGAP) interfaces.

\textbf{Key Responsibilities:}

\begin{itemize}
\tightlist
\item
  \textbf{NAS Termination:} All NAS messages (Registration, Service Request, PDU Session Establishment) terminate at AMF
\item
  \textbf{Registration Management:} Handles UE registration, deregistration, and periodic registration updates
\item
  \textbf{Connection Management:} Manages RRC signaling context via N2 toward gNB
\item
  \textbf{Mobility Management:} Handover signaling, tracking area management, paging coordination
\item
  \textbf{Security:} Initiates authentication (delegates to AUSF), manages NAS security context (ciphering + integrity)
\item
  \textbf{Slice Selection:} Interacts with NSSF to determine the appropriate network slice
\end{itemize}

\Needspace{16\baselineskip}

\textbf{Key Identifiers:}

{\def\LTcaptype{none} % do not increment counter
\Needspace{13\baselineskip}
\begin{longtable}[c]{@{}
  >{\raggedright\arraybackslash\hspace{0pt}}m{(\linewidth - 2\tabcolsep) * \real{0.5000}}
  >{\raggedright\arraybackslash\hspace{0pt}}m{(\linewidth - 2\tabcolsep) * \real{0.5000}}@{}}
\toprule\noalign{}
\rowcolor{CNink}\CNth{Identifier} & \CNth{Description} \\
\CNheadrulesep
\endhead
\bottomrule\noalign{}
\endlastfoot
\textbf{GUAMI} & Globally Unique AMF Identifier = MCC + MNC + AMF Region ID + AMF Set ID + AMF Pointer \\
\textbf{5G-GUTI} & 5G Globally Unique Temporary Identifier = GUAMI + 5G-TMSI (assigned to UE by AMF) \\
\textbf{AMF Region ID} & 8~bits — identifies a region within a PLMN \\
\textbf{AMF Set ID} & 10~bits — identifies an AMF Set within a region \\
\textbf{AMF Pointer} & 6~bits — identifies a specific AMF within a set \\
\end{longtable}
}

\textbf{Interfaces:}

\begin{itemize}
\tightlist
\item
  \textbf{N1:} UE \ensuremath{\leftrightarrow} AMF (NAS signaling, carried transparently over RAN)
\item
  \textbf{N2:} RAN \ensuremath{\leftrightarrow} AMF (NGAP protocol over SCTP)
\item
  \textbf{N8:} AMF \ensuremath{\leftrightarrow} UDM (subscription data retrieval)
\item
  \textbf{N11:} AMF \ensuremath{\leftrightarrow} SMF (session management messages)
\item
  \textbf{N12:} AMF \ensuremath{\leftrightarrow} AUSF (authentication procedures)
\item
  \textbf{N14:} AMF \ensuremath{\leftrightarrow} AMF (UE mobility between AMFs)
\item
  \textbf{N15:} AMF \ensuremath{\leftrightarrow} PCF (access and mobility policy)
\item
  \textbf{N22:} AMF \ensuremath{\leftrightarrow} NSSF (slice selection)
\end{itemize}

\textbf{Registration Area:}

\begin{itemize}
\tightlist
\item
  Composed of one or more Tracking Areas (TAs)
\item
  UE only performs Registration Update when moving outside its Registration Area
\item
  Larger registration areas \ensuremath{\to} less signaling, more paging load
\item
  AMF allocates the Registration Area based on UE mobility pattern
\end{itemize}

\CNsubsection{3.2}{Session Management Function (SMF)}

The SMF manages PDU sessions end-to-end. It selects and controls the UPF(s) that form the user plane path.

\textbf{Key Responsibilities:}

\begin{itemize}
\tightlist
\item
  \textbf{PDU Session Management:} Establishment, modification, and release of PDU sessions
\item
  \textbf{UPF Selection and Control:} Selects appropriate UPF(s), programs forwarding rules via N4/\allowbreak{}PFCP
\item
  \textbf{IP Address Allocation:} Assigns IPv4/\allowbreak{}IPv6 addresses to UEs (or delegates to external DHCP)
\item
  \textbf{QoS Enforcement:} Installs QoS rules in UPF based on PCF policies
\item
  \textbf{Roaming:} Handles home-routed and local breakout scenarios
\item
  \textbf{Downlink Data Notification:} Triggers paging via AMF when downlink data arrives for idle UE
\end{itemize}

\Needspace{16\baselineskip}

\textbf{PDU Session Types:}

{\def\LTcaptype{none} % do not increment counter
\Needspace{13\baselineskip}
\begin{longtable}[c]{@{}
  >{\raggedright\arraybackslash\hspace{0pt}}m{(\linewidth - 4\tabcolsep) * \real{0.3333}}
  >{\raggedright\arraybackslash\hspace{0pt}}m{(\linewidth - 4\tabcolsep) * \real{0.3333}}
  >{\raggedright\arraybackslash\hspace{0pt}}m{(\linewidth - 4\tabcolsep) * \real{0.3333}}@{}}
\toprule\noalign{}
\rowcolor{CNink}\CNth{Type} & \CNth{Description} & \CNth{Use Case} \\
\CNheadrulesep
\endhead
\bottomrule\noalign{}
\endlastfoot
\textbf{IPv4} & Traditional IPv4 connectivity & Legacy internet access \\
\textbf{IPv6} & IPv6-only connectivity & IoT, modern networks \\
\textbf{IPv4v6} & Dual-stack IPv4 + IPv6 & Default smartphone access \\
\textbf{Ethernet} & Layer-2 Ethernet frames & Industrial LAN extension, TSN \\
\textbf{Unstructured} & Raw IP or non-IP data & IoT sensors, point-to-point tunnels \\
\end{longtable}
}

\textbf{N4 Interface and PFCP:}

\begin{itemize}
\tightlist
\item
  \textbf{N4} connects SMF to UPF using \textbf{PFCP (Packet Forwarding Control Protocol)} — 3GPP TS 29.244
\item
  PFCP is a UDP-based protocol (port 8805) with:

  \begin{itemize}
  \tightlist
  \item
    \textbf{Session Establishment:} Create forwarding context in UPF
  \item
    \textbf{Session Modification:} Update rules (e.g., handover, QoS change)
  \item
    \textbf{Session Deletion:} Remove forwarding context
  \item
    \textbf{Heartbeat:} Detect UPF failure
  \end{itemize}
\end{itemize}

\textbf{Session Anchor:}

\begin{itemize}
\tightlist
\item
  The \textbf{PDU Session Anchor (PSA)} is the UPF that maintains a stable point of attachment to the Data Network
\item
  Provides IP address continuity during UE mobility
\item
  Does NOT change during handover (unlike intermediate UPFs)
\item
  SMF may insert Intermediate UPFs (I-UPF) for local routing while keeping PSA stable
\end{itemize}

\CNsubsection{3.3}{User Plane Function (UPF)}

The UPF is the \textbf{only user-plane NF} in the 5G Core. It handles all packet processing, forwarding, and policy enforcement for user data.

\textbf{Key Responsibilities:}

\begin{itemize}
\tightlist
\item
  \textbf{Packet Routing and Forwarding:} Routes packets between RAN and Data Network
\item
  \textbf{Traffic Detection and Reporting:} Performs DPI for application identification
\item
  \textbf{QoS Handling:} Enforces per-flow QoS (marking, policing, shaping)
\item
  \textbf{Usage Reporting:} Measures traffic volume/\allowbreak{}duration for charging
\item
  \textbf{Uplink Classifier (UL-CL):} Steers traffic to different data networks based on traffic matching
\item
  \textbf{Branching Point:} Supports multi-homed PDU sessions
\item
  \textbf{Buffering:} Buffers downlink packets for idle UEs, triggers DDN to SMF
\end{itemize}

\Needspace{14\baselineskip}

\textbf{Interfaces:}

{\def\LTcaptype{none} % do not increment counter
\Needspace{11\baselineskip}
\begin{longtable}[c]{@{}
  >{\raggedright\arraybackslash\hspace{0pt}}m{(\linewidth - 6\tabcolsep) * \real{0.2500}}
  >{\raggedright\arraybackslash\hspace{0pt}}m{(\linewidth - 6\tabcolsep) * \real{0.2500}}
  >{\raggedright\arraybackslash\hspace{0pt}}m{(\linewidth - 6\tabcolsep) * \real{0.2500}}
  >{\raggedright\arraybackslash\hspace{0pt}}m{(\linewidth - 6\tabcolsep) * \real{0.2500}}@{}}
\toprule\noalign{}
\rowcolor{CNink}\CNth{Interface} & \CNth{Connection} & \CNth{Protocol} & \CNth{Purpose} \\
\CNheadrulesep
\endhead
\bottomrule\noalign{}
\endlastfoot
\textbf{N3} & gNB \ensuremath{\leftrightarrow} UPF & GTP-U over UDP & User plane tunnel from RAN \\
\textbf{N6} & UPF \ensuremath{\leftrightarrow} DN & IP /\allowbreak{} Ethernet & Exit point to Data Network \\
\textbf{N9} & UPF \ensuremath{\leftrightarrow} UPF & GTP-U over UDP & Inter-UPF forwarding (chaining) \\
\textbf{N4} & SMF \ensuremath{\leftrightarrow} UPF & PFCP over UDP & Control of UPF from SMF \\
\end{longtable}
}

\Needspace{16\baselineskip}

\textbf{PFCP Rules (programmed by SMF via N4):}

{\def\LTcaptype{none} % do not increment counter
\Needspace{13\baselineskip}
\begin{longtable}[c]{@{}
  >{\raggedright\arraybackslash\hspace{0pt}}m{(\linewidth - 4\tabcolsep) * \real{0.3333}}
  >{\raggedright\arraybackslash\hspace{0pt}}m{(\linewidth - 4\tabcolsep) * \real{0.3333}}
  >{\raggedright\arraybackslash\hspace{0pt}}m{(\linewidth - 4\tabcolsep) * \real{0.3333}}@{}}
\toprule\noalign{}
\rowcolor{CNink}\CNth{Rule} & \CNth{Full Name} & \CNth{Purpose} \\
\CNheadrulesep
\endhead
\bottomrule\noalign{}
\endlastfoot
\textbf{PDR} & Packet Detection Rule & Matches incoming packets (source interface, IP filters, GTP TEID, QFI) \\
\textbf{FAR} & Forwarding Action Rule & Defines what to do (forward, drop, buffer, duplicate); specifies destination \\
\textbf{QER} & QoS Enforcement Rule & Rate limiting (MBR/\allowbreak{}GBR), gating (open/\allowbreak{}close), reflective QoS \\
\textbf{URR} & Usage Reporting Rule & Volume/\allowbreak{}time/\allowbreak{}event-based measurement triggers for charging \\
\textbf{BAR} & Buffering Action Rule & Controls buffering behavior, notification parameters \\
\end{longtable}
}

\textbf{Traffic Detection Methods:}

\begin{itemize}
\tightlist
\item
  \textbf{SDF Filter:} 5-tuple IP filter (src/\allowbreak{}dst IP, src/\allowbreak{}dst port, protocol)
\item
  \textbf{Application ID:} Pre-configured or dynamically learned via DPI
\item
  \textbf{GTP-U TEID:} Identifies specific GTP tunnels
\item
  \textbf{QFI (QoS Flow Identifier):} Maps to specific QoS flows
\end{itemize}

\CNsubsection{3.4}{Unified Data Management (UDM)}

The UDM provides subscriber data management, authentication credential handling, and subscription authorization. It is the 5G equivalent of the HSS.

\textbf{Key Responsibilities:}

\begin{itemize}
\tightlist
\item
  \textbf{SUPI Storage:} Stores and manages the Subscription Permanent Identifier (SUPI = IMSI or NAI format)
\item
  \textbf{Authentication Credential Computation:} Generates authentication vectors (AV) using subscriber’s long-term key (K) and operator parameters (OP/\allowbreak{}OPc)
\item
  \textbf{Subscription Profile Management:} Maintains subscribed services, DNN authorizations, QoS profiles
\item
  \textbf{Registration Management:} Tracks serving AMF/\allowbreak{}SMF per subscriber (NF registration)
\item
  \textbf{DNN Authorization:} Validates whether a UE is authorized to access a specific DNN/\allowbreak{}slice combination
\item
  \textbf{SUCI De-concealment:} Resolves SUCI back to SUPI (using SIDF — Subscription Identifier De-concealing Function)
\end{itemize}

\textbf{Key Data Stored:}

\begin{itemize}
\tightlist
\item
  SUPI, SUCI resolution keys
\item
  Subscriber authentication credentials (K, OP/\allowbreak{}OPc, SQN)
\item
  Subscribed S-NSSAI and DNN list
\item
  Default session QoS parameters
\item
  Serving NF registrations (AMF, SMF)
\item
  Access and mobility subscription data
\end{itemize}

\textbf{Interfaces:}

\begin{itemize}
\tightlist
\item
  \textbf{Nudm:} Service-based interface exposed by UDM
\item
  \textbf{N8 (via Nudm):} AMF retrieves subscription/\allowbreak{}mobility data
\item
  \textbf{N10 (via Nudm):} SMF retrieves session management subscription data
\item
  \textbf{N13 (via Nudm):} AUSF retrieves authentication credentials
\end{itemize}

\textbf{Relationship to UDR:}

\begin{itemize}
\tightlist
\item
  UDM contains application logic (auth vector computation, SUCI de-concealment)
\item
  UDR is the pure data store backend
\item
  UDM accesses UDR via \textbf{Nudr} interface
\end{itemize}

\CNsubsection{3.5}{Authentication Server Function (AUSF)}

The AUSF acts as the authentication server for the 5G network. It executes authentication procedures between the UE and the home network.

\textbf{Key Responsibilities:}

\begin{itemize}
\tightlist
\item
  \textbf{Authentication Execution:} Runs 5G-AKA or EAP-AKA’ procedures
\item
  \textbf{\ensuremath{K_{\mathrm{AUSF}}} Derivation:} Derives the anchor key (\ensuremath{K_{\mathrm{AUSF}}}) from which all subsequent keys are derived
\item
  \textbf{Authentication Result Storage:} Stores authentication status for future reference
\item
  \textbf{Home Network Authority:} Ensures authentication decision is made by the home network (even in roaming)
\end{itemize}

\Needspace{11\baselineskip}

\textbf{Authentication Methods:}

{\def\LTcaptype{none} % do not increment counter
\Needspace{8\baselineskip}
\begin{longtable}[c]{@{}
  >{\raggedright\arraybackslash\hspace{0pt}}m{(\linewidth - 4\tabcolsep) * \real{0.3333}}
  >{\raggedright\arraybackslash\hspace{0pt}}m{(\linewidth - 4\tabcolsep) * \real{0.3333}}
  >{\raggedright\arraybackslash\hspace{0pt}}m{(\linewidth - 4\tabcolsep) * \real{0.3333}}@{}}
\toprule\noalign{}
\rowcolor{CNink}\CNth{Method} & \CNth{Description} & \CNth{Use Case} \\
\CNheadrulesep
\endhead
\bottomrule\noalign{}
\endlastfoot
\textbf{5G-AKA} & 5G Authentication and Key Agreement — enhanced AKA with home network confirmation & Primary method for USIM-based UEs \\
\textbf{EAP-AKA’} & Extensible Authentication Protocol with AKA prime — EAP framework-based & Roaming, non-3GPP access, IoT \\
\end{longtable}
}

\CNkeepfig{m19-key-hierarchy}{2}

\textbf{Key Hierarchy from AUSF:}

\CNfigure{m19-key-hierarchy}{5G key hierarchy rooted at the AUSF}

\textbf{Interfaces:}

\begin{itemize}
\tightlist
\item
  \textbf{Nausf:} Service-based interface exposed by AUSF
\item
  \textbf{N12 (via Nausf):} AMF initiates authentication
\item
  \textbf{N13 (via Nudm):} AUSF retrieves authentication vectors from UDM
\end{itemize}

\CNsubsection{3.6}{Policy Control Function (PCF)}

The PCF is the centralized policy engine of the 5GC. It provides policy rules to control network behavior for access, sessions, and UE-specific policies.

\textbf{Key Responsibilities:}

\begin{itemize}
\tightlist
\item
  \textbf{PCC Rule Generation:} Creates Policy and Charging Control rules for QoS and charging
\item
  \textbf{AM Policy:} Access and Mobility policies (RFSP index, service area restrictions, RAT/\allowbreak{}frequency selection)
\item
  \textbf{SM Policy:} Session Management policies (QoS, gating, charging, traffic steering)
\item
  \textbf{UE Policy:} Policies delivered to the UE (URSP rules — UE Route Selection Policy, ANDSP)
\item
  \textbf{Policy Decision:} Makes real-time policy decisions based on subscription, network conditions, AF requests
\end{itemize}

\Needspace{12\baselineskip}

\textbf{Policy Types:}

{\def\LTcaptype{none} % do not increment counter
\Needspace{9\baselineskip}
\begin{longtable}[c]{@{}
  >{\raggedright\arraybackslash\hspace{0pt}}m{(\linewidth - 6\tabcolsep) * \real{0.2500}}
  >{\raggedright\arraybackslash\hspace{0pt}}m{(\linewidth - 6\tabcolsep) * \real{0.2500}}
  >{\raggedright\arraybackslash\hspace{0pt}}m{(\linewidth - 6\tabcolsep) * \real{0.2500}}
  >{\raggedright\arraybackslash\hspace{0pt}}m{(\linewidth - 6\tabcolsep) * \real{0.2500}}@{}}
\toprule\noalign{}
\rowcolor{CNink}\CNth{Policy} & \CNth{Scope} & \CNth{Delivered To} & \CNth{Purpose} \\
\CNheadrulesep
\endhead
\bottomrule\noalign{}
\endlastfoot
\textbf{AM Policy} & Access \& Mobility & AMF (N15) & Service area, RFSP, RAT restriction \\
\textbf{SM Policy} & Session & SMF (N7) & QoS flows, PCC rules, charging \\
\textbf{UE Policy} & UE behavior & UE (via AMF) & URSP, ANDSP, WLANSP \\
\end{longtable}
}

\textbf{PCC Rule Contents:}

\begin{itemize}
\tightlist
\item
  Rule identifier and precedence
\item
  Traffic flow template (SDF filter or application ID)
\item
  QoS parameters (5QI, ARP, MBR, GBR)
\item
  Gating status (open/\allowbreak{}close for UL/\allowbreak{}DL)
\item
  Charging method (online/\allowbreak{}offline) and rating group
\item
  Traffic steering policy
\end{itemize}

\textbf{Interfaces:}

\begin{itemize}
\tightlist
\item
  \textbf{Npcf:} Service-based interface exposed by PCF
\item
  \textbf{N7 (via Npcf):} SMF retrieves SM policy decisions
\item
  \textbf{N15 (via Npcf):} AMF retrieves AM policy decisions
\item
  \textbf{N5 (via Npcf):} AF influences policy via NEF or directly
\end{itemize}

\CNsubsection{3.7}{Network Slice Selection Function (NSSF)}

The NSSF assists in selecting the appropriate network slice instance and AMF set for a UE.

\textbf{Key Responsibilities:}

\begin{itemize}
\tightlist
\item
  \textbf{Allowed NSSAI Determination:} Determines which S-NSSAIs the UE is permitted to use in the serving PLMN
\item
  \textbf{S-NSSAI Mapping:} Maps Requested NSSAI to Configured NSSAI (home\ensuremath{\leftrightarrow}serving PLMN mapping in roaming)
\item
  \textbf{AMF Redirection:} When the current AMF cannot serve the selected slice, NSSF identifies a target AMF set that can
\item
  \textbf{NRF Selection:} Helps select appropriate NRF instance for slice-specific NF discovery
\end{itemize}

\Needspace{16\baselineskip}

\textbf{Key Concepts:}

{\def\LTcaptype{none} % do not increment counter
\Needspace{13\baselineskip}
\begin{longtable}[c]{@{}
  >{\raggedright\arraybackslash\hspace{0pt}}m{(\linewidth - 2\tabcolsep) * \real{0.5000}}
  >{\raggedright\arraybackslash\hspace{0pt}}m{(\linewidth - 2\tabcolsep) * \real{0.5000}}@{}}
\toprule\noalign{}
\rowcolor{CNink}\CNth{Term} & \CNth{Description} \\
\CNheadrulesep
\endhead
\bottomrule\noalign{}
\endlastfoot
\textbf{S-NSSAI} & Single NSSAI = SST (8-bit Slice/\allowbreak{}Service Type) + optional SD (24-bit Slice Differentiator) \\
\textbf{Requested NSSAI} & S-NSSAIs sent by UE in Registration Request \\
\textbf{Allowed NSSAI} & S-NSSAIs permitted by network (returned in Registration Accept) \\
\textbf{Configured NSSAI} & S-NSSAIs provisioned in the UE by the HPLMN/\allowbreak{}VPLMN \\
\textbf{Subscribed S-NSSAI} & Slice subscription in UDM \\
\end{longtable}
}

\textbf{Standard SST Values:}

\begin{itemize}
\tightlist
\item
  SST=1: eMBB (enhanced Mobile Broadband)
\item
  SST=2: URLLC (Ultra-Reliable Low-Latency Communications)
\item
  SST=3: MIoT (Massive IoT)
\item
  SST=4: V2X (Vehicle-to-Everything)
\end{itemize}

\CNsubsection{3.8}{Network Repository Function (NRF)}

The NRF is the central registry for all NF instances in the 5GC. It enables service discovery and authorization.

\textbf{Key Responsibilities:}

\begin{itemize}
\tightlist
\item
  \textbf{NF Registration:} NFs register their profiles (type, capacity, supported services, PLMN, slice, location)
\item
  \textbf{NF Discovery:} Consumer NFs query NRF to find producer NFs matching criteria
\item
  \textbf{NF Status Notification:} Notifies subscribed NFs about NF status changes (registered, deregistered, suspended)
\item
  \textbf{OAuth2 Token Issuance:} Issues access tokens for service-based interface authorization
\end{itemize}

\textbf{NF Profile Contents:}

\begin{itemize}
\tightlist
\item
  NF Instance ID (UUID)
\item
  NF Type (AMF, SMF, UPF, etc.)
\item
  NF Status (REGISTERED, SUSPENDED, UNDISCOVERABLE)
\item
  PLMN ID(s) served
\item
  S-NSSAI(s) supported
\item
  FQDN or IP addresses
\item
  Supported services and versions
\item
  Capacity, priority, load information
\item
  Geographic locality
\end{itemize}

\textbf{Discovery Query Parameters:}

\begin{itemize}
\tightlist
\item
  Target NF Type
\item
  Requester NF Type
\item
  Service names
\item
  Target PLMN, S-NSSAI, DNN
\item
  Preferred locality
\end{itemize}

\textbf{OAuth2 in 5GC:}

\begin{itemize}
\tightlist
\item
  NRF acts as OAuth2 Authorization Server
\item
  Consumer NF requests token with target NF type and services
\item
  Token included in SBI requests (Authorization header)
\item
  Producer NF validates token before serving the request
\end{itemize}

\CNsubsection{3.9}{Network Exposure Function (NEF)}

The NEF provides a secure framework for exposing 5GC capabilities and events to external Application Functions (AFs).

\textbf{Key Responsibilities:}

\begin{itemize}
\tightlist
\item
  \textbf{Exposure Framework:} Exposes network capabilities to 3rd-party AFs via standardized APIs
\item
  \textbf{Traffic Influence:} AF can request specific traffic routing (e.g., route traffic to local edge server)
\item
  \textbf{Monitoring Events:} AF subscribes to UE events (reachability, location, connectivity status)
\item
  \textbf{Parameter Provisioning:} AF provisions expected UE behavior, background data transfer policies
\item
  \textbf{Translation:} Maps between external identifiers (GPSI, external group ID) and internal identifiers (SUPI)
\item
  \textbf{Security Gateway:} Prevents direct AF access to internal NFs
\end{itemize}

\Needspace{17\baselineskip}

\textbf{Key Exposure Capabilities:}

{\def\LTcaptype{none} % do not increment counter
\Needspace{14\baselineskip}
\begin{longtable}[c]{@{}
  >{\raggedright\arraybackslash\hspace{0pt}}m{(\linewidth - 2\tabcolsep) * \real{0.5000}}
  >{\raggedright\arraybackslash\hspace{0pt}}m{(\linewidth - 2\tabcolsep) * \real{0.5000}}@{}}
\toprule\noalign{}
\rowcolor{CNink}\CNth{Capability} & \CNth{Description} \\
\CNheadrulesep
\endhead
\bottomrule\noalign{}
\endlastfoot
\textbf{Monitoring Events} & UE reachability, location reporting, loss of connectivity, roaming status \\
\textbf{Traffic Influence} & Request traffic routing to local Data Network (edge computing) \\
\textbf{Analytics Exposure} & Expose NWDAF analytics to external AFs \\
\textbf{Parameter Provisioning} & Expected UE behavior, background data transfer \\
\textbf{Device Triggering} & Trigger specific action on IoT device \\
\textbf{NIDD} & Non-IP Data Delivery for IoT \\
\end{longtable}
}

\CNsubsection{3.10}{Unified Data Repository (UDR)}

The UDR is the backend data store for subscriber, policy, and exposure data. It is accessed by UDM, PCF, and NEF.

\textbf{Key Responsibilities:}

\begin{itemize}
\tightlist
\item
  \textbf{Subscription Data Storage:} Stores subscriber profiles accessed by UDM
\item
  \textbf{Policy Data Storage:} Stores policy data (PCC rules, AM/\allowbreak{}SM policies) accessed by PCF
\item
  \textbf{Exposure Data Storage:} Stores exposure subscriptions and data accessed by NEF
\item
  \textbf{Application Data Storage:} Stores AF-provided data (e.g., PFDs — Packet Flow Descriptions)
\end{itemize}

\Needspace{14\baselineskip}

\textbf{Data Categories:}

{\def\LTcaptype{none} % do not increment counter
\Needspace{11\baselineskip}
\begin{longtable}[c]{@{}
  >{\raggedright\arraybackslash\hspace{0pt}}m{(\linewidth - 4\tabcolsep) * \real{0.3333}}
  >{\raggedright\arraybackslash\hspace{0pt}}m{(\linewidth - 4\tabcolsep) * \real{0.3333}}
  >{\raggedright\arraybackslash\hspace{0pt}}m{(\linewidth - 4\tabcolsep) * \real{0.3333}}@{}}
\toprule\noalign{}
\rowcolor{CNink}\CNth{Data Type} & \CNth{Consumer NF} & \CNth{Examples} \\
\CNheadrulesep
\endhead
\bottomrule\noalign{}
\endlastfoot
\textbf{Subscription Data} & UDM & Authentication credentials, subscribed S-NSSAI, DNN configs \\
\textbf{Policy Data} & PCF & PCC rule templates, sponsored data policies \\
\textbf{Exposure Data} & NEF & Monitoring subscriptions, AF influence requests \\
\textbf{Application Data} & NEF/\allowbreak{}PCF & PFDs, background data transfer policies \\
\end{longtable}
}

\textbf{Interface:}

\begin{itemize}
\tightlist
\item
  \textbf{Nudr:} Service-based interface (the only interface of UDR)
\item
  Accessed by UDM (Nudr\_\allowbreak{}DataManagement), PCF (Nudr\_\allowbreak{}DataManagement), NEF (Nudr\_\allowbreak{}DataManagement)
\end{itemize}

\CNsubsection{3.11}{Unstructured Data Storage Function (UDSF)}

The UDSF enables stateless design of NFs by providing external storage for NF context data.

\textbf{Key Responsibilities:}

\begin{itemize}
\tightlist
\item
  \textbf{Unstructured Data Storage:} Stores arbitrary NF context/\allowbreak{}state as opaque data blocks
\item
  \textbf{Stateless NF Support:} NFs can recover state after restart or load-balance across instances
\item
  \textbf{CRUD Operations:} Create, Read, Update, Delete operations on stored records
\end{itemize}

\textbf{Why UDSF Matters:}

\begin{itemize}
\tightlist
\item
  Without UDSF, each NF instance holds session state in memory
\item
  With UDSF, any NF instance can serve any request by retrieving state from UDSF
\item
  Enables true horizontal scaling, failover, and cloud-native deployment
\item
  Each NF type may have its own UDSF instance
\end{itemize}

\textbf{Interface:}

\begin{itemize}
\tightlist
\item
  \textbf{Nudsf:} Service-based interface for storing/\allowbreak{}retrieving unstructured data
\end{itemize}

\CNsubsection{3.12}{Service Communication Proxy (SCP)}

The SCP is a routing and load-balancing element placed in the SBI communication path between NFs. Defined in TS 23.501 Release 16+.

\textbf{Key Responsibilities:}

\begin{itemize}
\tightlist
\item
  \textbf{Indirect Communication:} NFs send requests to SCP instead of directly to the target NF
\item
  \textbf{Load Balancing:} Distributes requests across multiple NF instances
\item
  \textbf{Message Routing:} Routes messages based on discovery parameters without consumer NF needing full discovery
\item
  \textbf{Topology Hiding:} Hides internal NF topology from external networks (roaming)
\item
  \textbf{Overload Control:} Implements congestion handling and request throttling
\end{itemize}

\Needspace{14\baselineskip}

\textbf{Communication Models:}

{\def\LTcaptype{none} % do not increment counter
\Needspace{11\baselineskip}
\begin{longtable}[c]{@{}
  >{\raggedright\arraybackslash\hspace{0pt}}m{(\linewidth - 2\tabcolsep) * \real{0.5000}}
  >{\raggedright\arraybackslash\hspace{0pt}}m{(\linewidth - 2\tabcolsep) * \real{0.5000}}@{}}
\toprule\noalign{}
\rowcolor{CNink}\CNth{Model} & \CNth{Description} \\
\CNheadrulesep
\endhead
\bottomrule\noalign{}
\endlastfoot
\textbf{Model A (Direct)} & NF consumer discovers target via NRF, communicates directly \\
\textbf{Model B (Indirect w/\allowbreak{}o delegated discovery)} & NF discovers target via NRF, routes via SCP \\
\textbf{Model C (Indirect w/\allowbreak{} delegated discovery)} & NF sends to SCP, SCP performs discovery and routing \\
\textbf{Model D (Indirect w/\allowbreak{} delegated discovery + selection)} & SCP handles everything: discovery, selection, routing \\
\end{longtable}
}

\textbf{Benefits:}

\begin{itemize}
\tightlist
\item
  Simplifies NF implementation (NFs don’t need complex discovery/\allowbreak{}retry logic)
\item
  Centralized observability and telemetry
\item
  Easier deployment of cross-cutting concerns (mTLS termination, tracing)
\item
  Enables seamless NF scaling without reconfiguration
\end{itemize}

\CNsection{4}{Control Plane vs User Plane in 5G Core}

The 5GC enforces a strict separation between Control Plane (CP) and User Plane (UP), known as CUPS (Control and User Plane Separation).

\CNsubsection{}{Control Plane (CP)}

\begin{itemize}
\tightlist
\item
  \textbf{All NFs except UPF} operate in the control plane
\item
  Communication via \textbf{Service-Based Interfaces (SBI)} using HTTP/\allowbreak{}2 + JSON
\item
  Protocol stack: HTTP/\allowbreak{}2 \ensuremath{\to} TCP \ensuremath{\to} IP (or via SCP)
\item
  Functions: signaling, policy, authentication, registration, session management, discovery
\item
  Deployed centrally or regionally depending on latency requirements
\item
  AMF and SMF are the primary CP anchors for UE signaling and session management
\end{itemize}

\CNsubsection{}{User Plane (UP)}

\begin{itemize}
\tightlist
\item
  \textbf{Only UPF} operates in the user plane
\item
  Programmed by SMF via \textbf{N4 /\allowbreak{} PFCP} (control interface)
\item
  User data protocol stack: App Data \ensuremath{\to} PDU Layer \ensuremath{\to} GTP-U \ensuremath{\to} UDP \ensuremath{\to} IP
\item
  Functions: packet forwarding, QoS enforcement, traffic detection, usage measurement, buffering
\item
  Can be deployed at edge, regional, or central locations depending on latency and data locality needs
\item
  Multiple UPFs can be chained (via N9) for service function chaining or local breakout
\end{itemize}

\CNkeepfig{m19-cp-up}{3}

\CNsubsection{}{CP/\allowbreak{}UP Interaction Summary}

\CNfigure{m19-cp-up}{Control plane (SBI) and user plane of the 5G core}

\Needspace{28\baselineskip}

\CNsection{5}{Major Comparison: 4G EPC vs 5G Core}

{\def\LTcaptype{none} % do not increment counter
\Needspace{22\baselineskip}
\begin{longtable}[c]{@{}
  >{\raggedright\arraybackslash\hspace{0pt}}m{(\linewidth - 4\tabcolsep) * \real{0.2500}}
  >{\raggedright\arraybackslash\hspace{0pt}}m{(\linewidth - 4\tabcolsep) * \real{0.2500}}
  >{\raggedright\arraybackslash\hspace{0pt}}m{(\linewidth - 4\tabcolsep) * \real{0.5000}}@{}}
\toprule\noalign{}
\rowcolor{CNink}\CNth{Aspect} & \CNth{4G EPC} & \CNth{5G Core (5GC)} \\
\CNheadrulesep
\endhead
\bottomrule\noalign{}
\endlastfoot
\textbf{Architecture Style} & Point-to-point interfaces between monolithic NEs & Service-Based Architecture (SBA), microservices \\
\textbf{Signaling Protocol} & GTP-C (v2), Diameter & HTTP/\allowbreak{}2 + JSON (RESTful APIs) \\
\textbf{User Plane} & S-GW + P-GW (combined or separated) & UPF (single user-plane NF) \\
\textbf{Subscriber Database} & HSS (Diameter-based) & UDM + UDR (SBI-based, separated logic/\allowbreak{}storage) \\
\textbf{Authentication} & EPS-AKA via HSS & 5G-AKA /\allowbreak{} EAP-AKA’ via AUSF + UDM \\
\textbf{Policy} & PCRF (Gx, Rx, Sd interfaces) & PCF (Npcf SBI, AM/\allowbreak{}SM/\allowbreak{}UE policies) \\
\textbf{Session Management} & MME + S-GW + P-GW combined & SMF (dedicated NF, decoupled from mobility) \\
\textbf{Mobility Management} & MME (combined with session signaling) & AMF (pure mobility/\allowbreak{}access, no session logic) \\
\textbf{Network Slicing} & Not natively supported & Native support (S-NSSAI, NSSF, per-slice NF instances) \\
\textbf{Interface Style} & Fixed reference points (S1, S5, S6a, S11…) & SBI services + reference points for UP \\
\textbf{State Management} & Stateful NEs & Stateless NFs + UDSF for external state \\
\textbf{Deployment} & Appliance-based, monolithic & Cloud-native, containerized, microservices \\
\textbf{Discovery} & Preconfigured or DNS-based & Dynamic via NRF (registration + discovery + OAuth2) \\
\textbf{Edge Computing} & Difficult (LIPA/\allowbreak{}SIPTO limited) & Native (distributed UPF, LADN, MEC integration) \\
\end{longtable}
}

\CNkeepfig{m19-4g-to-5g}{3}

\CNsubsection{}{4G to 5G Network Function Mapping}

\CNfigure{m19-4g-to-5g}{How the EPC functions map to 5G core network functions}

\Needspace{33\baselineskip}

\CNsection{6}{Reference Points}

Reference points define the interfaces between specific NF pairs. While SBI is used for CP-to-CP communication, reference points remain important for understanding protocol specifics.

\CNsubsection{}{Complete Reference Point Table}

{\def\LTcaptype{none} % do not increment counter
\Needspace{22\baselineskip}
\begin{longtable}[c]{@{}
  >{\raggedright\arraybackslash\hspace{0pt}}m{(\linewidth - 6\tabcolsep) * \real{0.2500}}
  >{\raggedright\arraybackslash\hspace{0pt}}m{(\linewidth - 6\tabcolsep) * \real{0.2500}}
  >{\raggedright\arraybackslash\hspace{0pt}}m{(\linewidth - 6\tabcolsep) * \real{0.2500}}
  >{\raggedright\arraybackslash\hspace{0pt}}m{(\linewidth - 6\tabcolsep) * \real{0.2500}}@{}}
\toprule\noalign{}
\rowcolor{CNink}\CNth{Reference Point} & \CNth{Connection} & \CNth{Protocol/\allowbreak{}Transport} & \CNth{Purpose} \\
\CNheadrulesep
\endhead
\bottomrule\noalign{}
\endlastfoot
\textbf{N1} & UE \ensuremath{\leftrightarrow} AMF & NAS over RRC (via gNB) & NAS signaling (registration, session mgmt requests) \\
\textbf{N2} & RAN \ensuremath{\leftrightarrow} AMF & NGAP over SCTP & RAN-CN signaling (context setup, handover, paging) \\
\textbf{N3} & RAN \ensuremath{\leftrightarrow} UPF & GTP-U over UDP/\allowbreak{}IP & User plane tunnel carrying PDU session data \\
\textbf{N4} & SMF \ensuremath{\leftrightarrow} UPF & PFCP over UDP (port 8805) & UPF programming (PDR, FAR, QER, URR rules) \\
\textbf{N6} & UPF \ensuremath{\leftrightarrow} DN & IP /\allowbreak{} Ethernet & User traffic exit to Data Network \\
\textbf{N7} & SMF \ensuremath{\leftrightarrow} PCF & SBI (HTTP/\allowbreak{}2) & SM policy (PCC rules, QoS decisions) \\
\textbf{N8} & AMF \ensuremath{\leftrightarrow} UDM & SBI (HTTP/\allowbreak{}2) & Access/\allowbreak{}mobility subscription data \\
\textbf{N9} & UPF \ensuremath{\leftrightarrow} UPF & GTP-U over UDP/\allowbreak{}IP & Inter-UPF user plane forwarding \\
\textbf{N10} & SMF \ensuremath{\leftrightarrow} UDM & SBI (HTTP/\allowbreak{}2) & Session management subscription data \\
\textbf{N11} & AMF \ensuremath{\leftrightarrow} SMF & SBI (HTTP/\allowbreak{}2) & PDU session lifecycle management \\
\textbf{N12} & AMF \ensuremath{\leftrightarrow} AUSF & SBI (HTTP/\allowbreak{}2) & Authentication initiation \\
\textbf{N13} & AUSF \ensuremath{\leftrightarrow} UDM & SBI (HTTP/\allowbreak{}2) & Authentication vector retrieval \\
\textbf{N14} & AMF \ensuremath{\leftrightarrow} AMF & SBI (HTTP/\allowbreak{}2) & Inter-AMF mobility (UE context transfer) \\
\textbf{N15} & AMF \ensuremath{\leftrightarrow} PCF & SBI (HTTP/\allowbreak{}2) & AM policy (access/\allowbreak{}mobility restrictions) \\
\textbf{N22} & AMF \ensuremath{\leftrightarrow} NSSF & SBI (HTTP/\allowbreak{}2) & Slice selection assistance \\
\textbf{N27} & NRF \ensuremath{\leftrightarrow} NRF & SBI (HTTP/\allowbreak{}2) & Inter-PLMN NF discovery (via SEPP) \\
\textbf{N32} & SEPP \ensuremath{\leftrightarrow} SEPP & PRINS /\allowbreak{} TLS & Inter-PLMN security (roaming) \\
\end{longtable}
}

\CNsubsection{}{Key Reference Point Details}

\textbf{N1 (UE \ensuremath{\leftrightarrow} AMF):}

\begin{itemize}
\tightlist
\item
  Carries NAS PDUs (Registration Request/\allowbreak{}Accept, Authentication, Security Mode, PDU Session Establishment)
\item
  NAS messages are transparently relayed by gNB (encapsulated in NGAP on N2)
\item
  Supports NAS security (integrity + ciphering) after security activation
\end{itemize}

\textbf{N2 (RAN \ensuremath{\leftrightarrow} AMF):}

\begin{itemize}
\tightlist
\item
  NGAP protocol over SCTP
\item
  Procedures: Initial UE Message, PDU Session Resource Setup, Handover Required/\allowbreak{}Request, Paging
\item
  Carries UE-associated and non-UE-associated signaling
\end{itemize}

\textbf{N3 (RAN \ensuremath{\leftrightarrow} UPF):}

\begin{itemize}
\tightlist
\item
  GTP-U tunnel identified by TEID (Tunnel Endpoint Identifier)
\item
  One GTP-U tunnel per PDU session per gNB-UPF pair
\item
  QoS Flow Identifier (QFI) carried in GTP-U extension header to map QoS flows
\end{itemize}

\textbf{N4 (SMF \ensuremath{\leftrightarrow} UPF):}

\begin{itemize}
\tightlist
\item
  PFCP (Packet Forwarding Control Protocol) — 3GPP TS 29.244
\item
  Session-level management (Establishment, Modification, Deletion)
\item
  Node-level management (Association Setup, Heartbeat)
\item
  Reports: Usage Reports, Downlink Data Notification, End Marker
\end{itemize}

\textbf{N9 (UPF \ensuremath{\leftrightarrow} UPF):}

\begin{itemize}
\tightlist
\item
  GTP-U tunnel for multi-hop user plane paths
\item
  Used in UL-CL (Uplink Classifier) and branching point scenarios
\item
  Enables service function chaining and distributed UPF deployment
\end{itemize}

\CNsection{7}{Deployment Options}

The 5GC’s modular architecture allows flexible deployment topologies based on service requirements.

\CNsubsection{7.1}{Centralized Deployment}

All NFs (including UPF) deployed in a central data center.

\textbf{Characteristics:}

\begin{itemize}
\tightlist
\item
  Simple operations, single site
\item
  Higher latency for user plane (all traffic traverses central site)
\item
  Suitable for wide-area coverage with relaxed latency requirements
\item
  Lower infrastructure cost
\end{itemize}

\CNsubsection{7.2}{Distributed Deployment}

Control plane centralized, user plane distributed closer to users.

\textbf{Characteristics:}

\begin{itemize}
\tightlist
\item
  CP NFs (AMF, SMF, PCF, etc.) in regional/\allowbreak{}central DC
\item
  UPF deployed at regional aggregation points
\item
  Reduced user-plane latency
\item
  SMF manages multiple UPF instances across locations
\end{itemize}

\CNsubsection{7.3}{Edge UPF Deployment (MEC Integration)}

UPF instances deployed at the network edge (cell site or local aggregation point) for ultra-low latency.

\textbf{Characteristics:}

\begin{itemize}
\tightlist
\item
  Local breakout of traffic at the edge
\item
  Enables Multi-access Edge Computing (MEC) — ETSI MEC framework
\item
  SMF can insert/\allowbreak{}remove edge UPFs dynamically based on UE location
\item
  \textbf{UL-CL (Uplink Classifier):} Steers specific traffic flows to edge UPF while rest goes to central UPF
\item
  \textbf{Local Area Data Network (LADN):} Data Network available only in specific geographic area
\item
  Typical use cases: AR/\allowbreak{}VR, autonomous driving, industrial automation, cloud gaming
\end{itemize}

\Needspace{18\baselineskip}

\CNsubsection{}{Deployment Architecture Comparison}

{\def\LTcaptype{none} % do not increment counter
\Needspace{14\baselineskip}
\begin{longtable}[c]{@{}
  >{\raggedright\arraybackslash\hspace{0pt}}m{(\linewidth - 6\tabcolsep) * \real{0.1860}}
  >{\raggedright\arraybackslash\hspace{0pt}}m{(\linewidth - 6\tabcolsep) * \real{0.2791}}
  >{\raggedright\arraybackslash\hspace{0pt}}m{(\linewidth - 6\tabcolsep) * \real{0.3023}}
  >{\raggedright\arraybackslash\hspace{0pt}}m{(\linewidth - 6\tabcolsep) * \real{0.2326}}@{}}
\toprule\noalign{}
\rowcolor{CNink}\CNth{Aspect} & \CNth{Centralized} & \CNth{Distributed} & \CNth{Edge UPF} \\
\CNheadrulesep
\endhead
\bottomrule\noalign{}
\endlastfoot
\textbf{CP Location} & Central DC & Central/\allowbreak{}Regional DC & Central/\allowbreak{}Regional DC \\
\textbf{UP Location} & Central DC & Regional PoP & Cell site /\allowbreak{} Edge DC \\
\textbf{User Plane Latency} & High (10-50ms) & Medium (5-20ms) & Ultra-low (1-5ms) \\
\textbf{Complexity} & Low & Medium & High \\
\textbf{Use Cases} & General broadband & Regional services & URLLC, MEC, V2X \\
\textbf{UPF Count} & Few & Moderate & Many \\
\end{longtable}
}

\CNkeepfig{m19-pdu-path}{3}

\CNsubsection{}{PDU Session User Plane Path}

\CNfigure{m19-pdu-path}{User-plane path of a PDU session with an intermediate UPF}

\textbf{PDU Session Establishment Flow (simplified):}

\begin{enumerate}
\def\labelenumi{\arabic{enumi}.}
\tightlist
\item
  UE \ensuremath{\to} AMF: NAS PDU Session Establishment Request (via N1)
\item
  AMF \ensuremath{\to} SMF: Nsmf\_\allowbreak{}PDUSession\_\allowbreak{}CreateSMContext (via N11)
\item
  SMF selects UPF, allocates IP address
\item
  SMF \ensuremath{\to} UPF: PFCP Session Establishment (via N4) — installs PDR/\allowbreak{}FAR/\allowbreak{}QER
\item
  SMF \ensuremath{\to} AMF: N1N2MessageTransfer (contains NAS accept + N2 info)
\item
  AMF \ensuremath{\to} gNB: PDU Session Resource Setup Request (via N2)
\item
  gNB \ensuremath{\to} UPF: GTP-U tunnel established (N3)
\item
  UE receives PDU Session Establishment Accept with QoS rules
\end{enumerate}

\CNsection{}{Summary}

The 5G Core represents a paradigm shift from the 4G EPC:

\begin{itemize}
\tightlist
\item
  \textbf{SBA replaces point-to-point interfaces} — enabling flexibility, reusability, and independent NF evolution
\item
  \textbf{Complete CUPS} — UPF is the only user-plane entity, programmable via PFCP
\item
  \textbf{Cloud-native by design} — stateless NFs, UDSF, NRF-based discovery, containerization
\item
  \textbf{Native slicing} — S-NSSAI, NSSF, per-slice NF selection from the ground up
\item
  \textbf{Security enhanced} — SUCI concealment, home-network-controlled authentication, SBI OAuth2
\item
  \textbf{Edge-ready} — distributed UPF, UL-CL, LADN for ultra-low-latency services
\end{itemize}

\emph{References: 3GPP TS 23.501 (System Architecture), TS 23.502 (Procedures), TS 29.244 (PFCP), TS 33.501 (Security Architecture)}
